\documentclass[uplatex,dvipdfmx,11pt]{jsarticle}

\usepackage{amsmath,amssymb,amsthm,mathtools}
\usepackage[margin=25mm]{geometry}

\newcommand{\R}{\mathbf{R}}
\newcommand{\Z}{\mathbf{Z}}

\theoremstyle{definition}
\newtheorem{prop}{命題}
\newtheorem{remark}{注意}

\title{$T^2=\R^2/\Z^2$ の de Rham コホモロジー\\
--- 記号の省略をしない版 ---}
\author{}
\date{}

\begin{document}
\maketitle

\section{目的}

トーラス
\[
T^2=\R^2/\Z^2
\]
の de Rham コホモロジーを、被覆空間 \(\R^2\) 上の周期形式を用いて計算する。

ここでは記号の混乱を避けるため、
\begin{itemize}
  \item \(\R^2\) 上の標準座標に付随する $1$-形式を \(dx,dy\) と書く。
  \item それらが \(T^2\) に降りてできる $1$-形式を \(\alpha,\beta\) と書く。
  \item \(dx\wedge dy\) が \(T^2\) に降りてできる $2$-形式を \(\mu\) と書く。
\end{itemize}

最終的に
\[
H^0_{\mathrm{dR}}(T^2)\cong \R,\qquad
H^1_{\mathrm{dR}}(T^2)\cong \R^2,\qquad
H^2_{\mathrm{dR}}(T^2)\cong \R
\]
を示し、基底を
\[
[1],\qquad [\alpha],[\beta],\qquad [\mu]
\]
と記述する。

\section{基本設定}

商写像を
\[
\pi:\R^2\to T^2,\qquad \pi(x,y)=[(x,y)]
\]
とする。

\(\R^2\) 上の標準座標関数
\[
x:\R^2\to\R,\qquad (x,y)\mapsto x,
\]
\[
y:\R^2\to\R,\qquad (x,y)\mapsto y
\]
に対応する $1$-形式を
\[
dx=d(x),\qquad dy=d(y)
\]
と書く。

各 \((m,n)\in\Z^2\) に対し、整数平行移動
\[
\tau_{m,n}:\R^2\to\R^2,\qquad \tau_{m,n}(x,y)=(x+m,y+n)
\]
を考えると、
\[
\tau_{m,n}^*(dx)=d(x\circ \tau_{m,n})=d(x+m)=dx,
\]
\[
\tau_{m,n}^*(dy)=d(y\circ \tau_{m,n})=d(y+n)=dy
\]
である。したがって \(dx,dy\) は \(\Z^2\)-不変である。

また
\[
\tau_{m,n}^*(dx\wedge dy)=\tau_{m,n}^*(dx)\wedge \tau_{m,n}^*(dy)=dx\wedge dy
\]
であるから、\(dx\wedge dy\) も \(\Z^2\)-不変である。

\subsection{\(T^2\) 上の形式と \(\R^2\) 上の周期形式}

\begin{prop}[商写像は被覆写像]
\(\pi:\R^2\to T^2\) は滑らかな被覆写像である。とくに \(\pi\) は全射であり、
各点 \(p\in \R^2\) において微分
\[
d\pi_p:T_p\R^2\to T_{\pi(p)}T^2
\]
は線形同型である。
\end{prop}

\begin{proof}[略証]
\(\R^2/\Z^2\) の各点は、\(\R^2\) の小さい開矩形の像で局所的に記述でき、
その上で \(\pi\) は微分同相になる。したがって \(\pi\) は被覆写像であり、
とくに局所微分同相である。
\end{proof}

\begin{prop}[pullback の単射性]
各 \(k\ge 0\) に対して
\[
\pi^*:\Omega^k(T^2)\to \Omega^k(\R^2)
\]
は単射である。
\end{prop}

\begin{proof}
\(\eta\in\Omega^k(T^2)\) が \(\pi^*\eta=0\) を満たすとする。
任意の \(q\in T^2\) をとり、\(\pi(p)=q\) となる \(p\in\R^2\) をとる。
\(\pi\) は局所微分同相なので
\[
d\pi_p:T_p\R^2\to T_qT^2
\]
は同型である。したがって、任意の \(v_1,\dots,v_k\in T_qT^2\) に対し
\(d\pi_p(w_i)=v_i\) を満たす \(w_i\in T_p\R^2\) をとれる。

pullback の定義より
\[
(\pi^*\eta)_p(w_1,\dots,w_k)
=
\eta_q(d\pi_p(w_1),\dots,d\pi_p(w_k))
=
\eta_q(v_1,\dots,v_k).
\]
左辺は仮定より \(0\) なので
\[
\eta_q(v_1,\dots,v_k)=0.
\]
\(q\) と \(v_i\) は任意だったから \(\eta=0\) である。
\end{proof}

\begin{prop}[\(\Z^2\)-不変形式との対応]
\(\pi^*\) は \(T^2\) 上の微分形式と、\(\R^2\) 上の \(\Z^2\)-不変微分形式との間の
全単射を与える。すなわち
\[
\Omega^k(T^2)
\cong
\{\omega\in \Omega^k(\R^2)\mid \tau_{m,n}^*\omega=\omega\ \ (\forall (m,n)\in\Z^2)\}.
\]
\end{prop}

\begin{proof}
まず \(\eta\in\Omega^k(T^2)\) に対して
\[
\tau_{m,n}^*(\pi^*\eta)
=
(\pi\circ\tau_{m,n})^*\eta
=
\pi^*\eta
\]
である。したがって \(\pi^*\eta\) は \(\Z^2\)-不変である。

逆に \(\omega\in\Omega^k(\R^2)\) が \(\Z^2\)-不変であるとする。
各 \(q\in T^2\) と \(p\in\pi^{-1}(q)\) に対して
\[
\eta_q(v_1,\dots,v_k)
:=
\omega_p(w_1,\dots,w_k)
\]
と定める。ただし \(w_i\in T_p\R^2\) は
\[
d\pi_p(w_i)=v_i
\]
を満たすものとする。\(d\pi_p\) は同型なので \(w_i\) は一意に定まり、
さらに \(\omega\) の \(\Z^2\)-不変性により \(p\) の取り方にも依らない。
こうして \(\eta\in\Omega^k(T^2)\) が定まり、\(\pi^*\eta=\omega\) を満たす。
単射性は前命題より従う。
\end{proof}

\section{\(T^2\) 上の形式 \(\alpha,\beta,\mu\)}

前節の命題より、\(\Z^2\)-不変な形式 \(dx,dy,dx\wedge dy\) は
\(T^2\) 上の形式に一意に降りる。
それらを
\[
\alpha,\beta\in \Omega^1(T^2),\qquad \mu\in \Omega^2(T^2)
\]
と書く。すなわち
\[
\pi^*\alpha=dx,\qquad \pi^*\beta=dy,\qquad \pi^*\mu=dx\wedge dy
\tag{1}
\]
を満たすものとして定める。

\begin{remark}
ふつうの本では \(\alpha,\beta,\mu\) も単に \(dx,dy,dx\wedge dy\) と書くことが多いが、
本稿ではこの省略をしない。
\end{remark}

\section{\(H^0_{\mathrm{dR}}(T^2)\) の計算}

\begin{prop}
\[
H^0_{\mathrm{dR}}(T^2)\cong \R.
\]
\end{prop}

\begin{proof}
\(f\in \Omega^0(T^2)\) が閉であるとは \(df=0\) であることをいう。
このとき
\[
d(\pi^*f)=\pi^*(df)=0
\]
であるから、\(\pi^*f\) は \(\R^2\) 上局所定数である。
\(\R^2\) は連結なので、\(\pi^*f\) は定数関数である。
したがって \(f\) も定数関数である。

ゆえに
\[
H^0_{\mathrm{dR}}(T^2)\cong \R.
\]
\end{proof}

\section{\(H^1_{\mathrm{dR}}(T^2)\) の計算}

\subsection{2本の基本ループ}

\(T^2\) 上の path
\[
\gamma_x:[0,1]\to T^2,\qquad \gamma_x(t)=\pi(t,0),
\]
\[
\gamma_y:[0,1]\to T^2,\qquad \gamma_y(t)=\pi(0,t)
\]
を考える。これらは \(T^2\) 上では loop である。

それぞれの lift を
\[
\widetilde{\gamma}_x:[0,1]\to \R^2,\qquad \widetilde{\gamma}_x(t)=(t,0),
\]
\[
\widetilde{\gamma}_y:[0,1]\to \R^2,\qquad \widetilde{\gamma}_y(t)=(0,t)
\]
とする。すると
\[
\pi\circ \widetilde{\gamma}_x=\gamma_x,\qquad
\pi\circ \widetilde{\gamma}_y=\gamma_y.
\tag{2}
\]

\subsection{\(\alpha,\beta\) の基本ループに沿う積分}

\begin{prop}
\[
\int_{\gamma_x}\alpha=1,\qquad \int_{\gamma_x}\beta=0,
\]
\[
\int_{\gamma_y}\alpha=0,\qquad \int_{\gamma_y}\beta=1.
\]
\end{prop}

\begin{proof}
pullback の合成則 \((f\circ g)^*=g^*\circ f^*\) と \((1),(2)\) から
\[
\gamma_x^*\alpha=(\pi\circ \widetilde{\gamma}_x)^*\alpha
=\widetilde{\gamma}_x^{\,*}(\pi^*\alpha)
=\widetilde{\gamma}_x^{\,*}(dx),
\]
\[
\gamma_x^*\beta=(\pi\circ \widetilde{\gamma}_x)^*\beta
=\widetilde{\gamma}_x^{\,*}(\pi^*\beta)
=\widetilde{\gamma}_x^{\,*}(dy).
\]
さらに
\[
x\circ \widetilde{\gamma}_x(t)=t,\qquad y\circ \widetilde{\gamma}_x(t)=0
\]
より
\[
\widetilde{\gamma}_x^{\,*}(dx)=d(x\circ \widetilde{\gamma}_x)=d(t)=dt,
\]
\[
\widetilde{\gamma}_x^{\,*}(dy)=d(y\circ \widetilde{\gamma}_x)=d(0)=0.
\]
したがって
\[
\int_{\gamma_x}\alpha=\int_0^1 \gamma_x^*\alpha=\int_0^1 dt=1,
\]
\[
\int_{\gamma_x}\beta=\int_0^1 \gamma_x^*\beta=\int_0^1 0=0.
\]

同様に
\[
\gamma_y^*\alpha=(\pi\circ \widetilde{\gamma}_y)^*\alpha
=\widetilde{\gamma}_y^{\,*}(\pi^*\alpha)
=\widetilde{\gamma}_y^{\,*}(dx),
\]
\[
\gamma_y^*\beta=(\pi\circ \widetilde{\gamma}_y)^*\beta
=\widetilde{\gamma}_y^{\,*}(\pi^*\beta)
=\widetilde{\gamma}_y^{\,*}(dy).
\]
ここで
\[
x\circ \widetilde{\gamma}_y(t)=0,\qquad y\circ \widetilde{\gamma}_y(t)=t
\]
だから
\[
\widetilde{\gamma}_y^{\,*}(dx)=d(0)=0,\qquad
\widetilde{\gamma}_y^{\,*}(dy)=d(t)=dt.
\]
ゆえに
\[
\int_{\gamma_y}\alpha=\int_0^1 0=0,\qquad
\int_{\gamma_y}\beta=\int_0^1 dt=1.
\]
\end{proof}

\subsection{周期写像}

閉 $1$-形式 \(\omega\in\Omega^1(T^2)\) に対して
\[
\Phi([\omega])
=
\left(
\int_{\gamma_x}\omega,\,
\int_{\gamma_y}\omega
\right)\in \R^2
\]
と定める。

\(\omega=df\) が exact なら、任意の閉曲線 \(\gamma\) に対して
\[
\int_\gamma \omega=\int_\gamma df=0
\]
であるから、\(\Phi\) は well-defined であり、
\[
\Phi:H^1_{\mathrm{dR}}(T^2)\to \R^2
\]
という線形写像を与える。

\subsection{\(\Phi\) の全射性}

\begin{prop}
\(\Phi\) は全射である。
\end{prop}

\begin{proof}
前節の計算より
\[
\Phi([\alpha])=(1,0),\qquad \Phi([\beta])=(0,1).
\]
したがって \(\Phi\) は全射である。
\end{proof}

\begin{remark}[証明の見通し]
\(\Phi\) の単射性は \(\ker \Phi=0\) を示せばよい。そこで
\[
\Phi([\omega])=(0,0)
\]
を満たす閉 $1$-形式 \(\omega\) をとる。\(\R^2\) は可縮なので、\(\omega\) を
\(\R^2\) に引き戻した \(\pi^*\omega\) は exact であり、
\[
dg=\pi^*\omega
\]
を満たす \(g\in C^\infty(\R^2)\) が存在する。したがって残る問題は、この
\(g\) が \(\Z^2\)-周期的であり、よって \(T^2\) 上の関数に降りることを示す点に尽きる。

\medskip

実際、\(\Phi([\omega])=(0,0)\) という仮定は、まさにこの周期性を示すために使われる。
以下の証明ではこれを座標計算で示し、後の補足ではデッキ変換を用いた
より概念的な説明を与える。
\end{remark}


\subsection{\(\Phi\) の単射性}

\begin{prop}
\(\Phi\) は単射である。
\end{prop}

\begin{proof}
\(\Phi([\omega])=(0,0)\) と仮定する。
\(\omega\in\Omega^1(T^2)\) は閉であるとする。

\(\pi^*\omega\in \Omega^1(\R^2)\) は \(\Z^2\)-不変な $1$-形式であり、
ある滑らかな周期関数 \(A,B\in C^\infty(\R^2)\) を用いて
\[
\pi^*\omega=A(x,y)\,dx+B(x,y)\,dy
\tag{3}
\]
と書ける。

\(\omega\) が閉であるから
\[
0=\pi^*(d\omega)=d(\pi^*\omega).
\]
よって \((3)\) を微分して
\[
\frac{\partial A}{\partial y}=\frac{\partial B}{\partial x}
\tag{4}
\]
を得る。

\medskip

\(\R^2\) は可縮であるから、Poincar\'e の補題より \(\pi^*\omega\) は exact である。
したがって、ある \(g\in C^\infty(\R^2)\) が存在して
\[
dg=\pi^*\omega
\tag{5}
\]
を満たす。

したがって、あとは \(g\) が \(\Z^2\)-周期的であることを示せばよい。
ここで \(\Phi([\omega])=(0,0)\) という仮定は、その周期性を示すためにのみ用いる。

\medskip

まず
\[
I(y):=\int_0^1 A(x,y)\,dx
\]
とおく。
すると \((4)\) より
\[
I'(y)
=
\int_0^1 \frac{\partial A}{\partial y}(x,y)\,dx
=
\int_0^1 \frac{\partial B}{\partial x}(x,y)\,dx
=
B(1,y)-B(0,y)=0
\]
である。最後は \(B\) の \(1\)-periodicity を用いた。
したがって \(I(y)\) は定数である。

一方、
\[
\int_{\gamma_x}\omega=0
\]
と仮定しているので、\((2)\) と \((3)\) を用いて
\[
0=\int_{\gamma_x}\omega
=\int_0^1 \gamma_x^*\omega
=\int_0^1 \widetilde{\gamma}_x^{\,*}(\pi^*\omega)
=\int_0^1 \widetilde{\gamma}_x^{\,*}(A\,dx+B\,dy).
\]
ここで \(\widetilde{\gamma}_x(t)=(t,0)\) であるから
\[
\widetilde{\gamma}_x^{\,*}(A\,dx+B\,dy)=A(t,0)\,dt.
\]
よって
\[
0=\int_0^1 A(t,0)\,dt=I(0).
\]
\(I(y)\) は定数なので
\[
I(y)=0\qquad (\forall y).
\tag{6}
\]

次に \((5)\) より
\[
\frac{\partial g}{\partial x}=A.
\]
したがって
\[
g(x+1,y)-g(x,y)
=
\int_x^{x+1}\frac{\partial g}{\partial x}(s,y)\,ds
=
\int_x^{x+1}A(s,y)\,ds.
\]
\(A\) は \(1\)-periodic なので
\[
\int_x^{x+1}A(s,y)\,ds=\int_0^1A(s,y)\,ds=I(y)=0.
\]
ゆえに
\[
g(x+1,y)=g(x,y).
\tag{7}
\]

\medskip

同様に
\[
J(x):=\int_0^1 B(x,y)\,dy
\]
とおくと
\[
J'(x)
=
\int_0^1 \frac{\partial B}{\partial x}(x,y)\,dy
=
\int_0^1 \frac{\partial A}{\partial y}(x,y)\,dy
=
A(x,1)-A(x,0)=0
\]
より \(J(x)\) は定数である。

また \(\int_{\gamma_y}\omega=0\) だから
\[
0=\int_{\gamma_y}\omega
=\int_0^1 \gamma_y^*\omega
=\int_0^1 \widetilde{\gamma}_y^{\,*}(\pi^*\omega)
=\int_0^1 \widetilde{\gamma}_y^{\,*}(A\,dx+B\,dy).
\]
ここで \(\widetilde{\gamma}_y(t)=(0,t)\) だから
\[
\widetilde{\gamma}_y^{\,*}(A\,dx+B\,dy)=B(0,t)\,dt.
\]
よって
\[
0=\int_0^1 B(0,t)\,dt=J(0),
\]
したがって
\[
J(x)=0\qquad (\forall x).
\tag{8}
\]

さらに \((5)\) より
\[
\frac{\partial g}{\partial y}=B.
\]
したがって
\[
g(x,y+1)-g(x,y)
=
\int_y^{y+1}\frac{\partial g}{\partial y}(x,t)\,dt
=
\int_y^{y+1}B(x,t)\,dt
=
\int_0^1B(x,t)\,dt
=
J(x)=0.
\]
ゆえに
\[
g(x,y+1)=g(x,y).
\tag{9}
\]

\medskip

\((7),(9)\) より \(g\) は \(\Z^2\)-周期的である。
したがって、ある \(f\in C^\infty(T^2)\) が存在して
\[
\pi^*f=g
\]
を満たす。

すると
\[
\pi^*(df)=d(\pi^*f)=dg=\pi^*\omega.
\]
\(\pi^*\) の単射性より
\[
df=\omega.
\]
よって \(\omega\) は exact である。

したがって \(\ker\Phi=0\) であり、\(\Phi\) は単射である。
\end{proof}

\subsection{\(H^1\) の結論}

\begin{prop}
\[
H^1_{\mathrm{dR}}(T^2)\cong \R^2.
\]
しかも \([\alpha],[\beta]\) はその基底である。
\end{prop}

\begin{proof}
前節までで \(\Phi:H^1_{\mathrm{dR}}(T^2)\to\R^2\) が線形同型であることを示した。
また
\[
\Phi([\alpha])=(1,0),\qquad \Phi([\beta])=(0,1)
\]
であるから、\([\alpha],[\beta]\) は基底を与える。
\end{proof}

\begin{remark}
この計算は、閉 $1$-形式 \(\omega\) のコホモロジー類が
\[
\int_{\gamma_x}\omega,\qquad \int_{\gamma_y}\omega
\]
という2つの周期で完全に決まることを示している。
したがって \(H^1(T^2)\) は、横方向・縦方向という2つの独立な「回り方」を記録する空間だと理解できる。
\end{remark}

\section{\(H^2_{\mathrm{dR}}(T^2)\) の計算}

\subsection{積分写像}

閉 $2$-形式 \(\sigma\in \Omega^2(T^2)\) に対して
\[
I([\sigma])=\int_{T^2}\sigma
\]
と定める。

\(\sigma=d\lambda\) が exact なら、Stokes の定理より
\[
\int_{T^2}\sigma=\int_{T^2}d\lambda=\int_{\partial T^2}\lambda=0
\]
である。したがって
\[
I:H^2_{\mathrm{dR}}(T^2)\to \R
\]
は well-defined な線形写像である。

\subsection{\(I\) の全射性}

\begin{prop}
\(I\) は全射である。
\end{prop}

\begin{proof}
\((1)\) より \(\pi^*\mu=dx\wedge dy\) である。
したがって
\[
\int_{T^2}\mu=\int_{[0,1]^2} dx\,dy=1.
\]
ゆえに
\[
I([\mu])=1
\]
であり、\(I\) は全射である。
\end{proof}

\subsection{\(I\) の単射性}

\begin{prop}
\(I\) は単射である。
\end{prop}

\begin{proof}
\(I([\sigma])=0\) と仮定する。
\(\sigma\in \Omega^2(T^2)\) は閉 $2$-形式であるが、\(T^2\) は $2$ 次元多様体なので、
任意の $2$-形式は自動的に閉である。

\(\pi^*\sigma\in \Omega^2(\R^2)\) は \(\Z^2\)-不変な $2$-形式であり、
ある周期関数 \(h\in C^\infty(\R^2)\) を用いて
\[
\pi^*\sigma=h(x,y)\,dx\wedge dy
\tag{10}
\]
と書ける。

仮定 \(I([\sigma])=0\) は
\[
\int_{T^2}\sigma=0
\]
であるから、
\[
0=\int_{T^2}\sigma=\int_{[0,1]^2} h(x,y)\,dx\,dy
\tag{11}
\]
となる。

まず
\[
\phi(y):=\int_0^1 h(s,y)\,ds
\]
とおく。すると \(\phi\) は \(1\)-periodic であり、
\[
\int_0^1 \phi(y)\,dy
=
\int_0^1\int_0^1 h(s,y)\,ds\,dy
=0
\tag{12}
\]
が成り立つ。

次に
\[
u(x,y):=\int_0^x h(s,y)\,ds - x\,\phi(y)
\]
とおくと
\[
\frac{\partial u}{\partial x}(x,y)=h(x,y)-\phi(y).
\tag{13}
\]
また \(h,\phi\) の周期性から \(u\) は \(x,y\) の両方について \(1\)-periodic である。

さらに
\[
v(y):=\int_0^y \phi(t)\,dt
\]
とおくと
\[
v'(y)=\phi(y)
\tag{14}
\]
であり、\((12)\) から \(v\) も \(1\)-periodic である。

ここで \(\R^2\) 上の $1$-形式
\[
\widetilde{\lambda}:=-v(y)\,dx+u(x,y)\,dy
\]
を考える。係数 \(u,v\) は周期的であるから、\(\widetilde{\lambda}\) は \(\Z^2\)-不変である。
したがって、ある \(\lambda\in \Omega^1(T^2)\) が存在して
\[
\pi^*\lambda=\widetilde{\lambda}
\tag{15}
\]
を満たす。

\((15)\) を微分すると
\[
\pi^*(d\lambda)=d(\pi^*\lambda)=d\widetilde{\lambda}.
\]
しかも
\[
d\widetilde{\lambda}
=
d(-v(y)\,dx)+d(u(x,y)\,dy)
\]
\[
=
-v'(y)\,dy\wedge dx+\frac{\partial u}{\partial x}(x,y)\,dx\wedge dy
\]
\[
=
\left(v'(y)+\frac{\partial u}{\partial x}(x,y)\right)\,dx\wedge dy.
\]
ここに \((13),(14)\) を代入すると
\[
d\widetilde{\lambda}
=
\left(\phi(y)+h(x,y)-\phi(y)\right)\,dx\wedge dy
=
h(x,y)\,dx\wedge dy.
\]
\((10)\) より
\[
d\widetilde{\lambda}=\pi^*\sigma.
\]
したがって
\[
\pi^*(d\lambda)=\pi^*\sigma.
\]
\(\pi^*\) の単射性より
\[
d\lambda=\sigma.
\]
よって \(\sigma\) は exact である。

したがって \(\ker I=0\) であり、\(I\) は単射である。
\end{proof}

\subsection{\(H^2\) の結論}

\begin{prop}
\[
H^2_{\mathrm{dR}}(T^2)\cong \R.
\]
しかも \([\mu]\) はその基底である。
\end{prop}

\begin{proof}
前節までで \(I:H^2_{\mathrm{dR}}(T^2)\to\R\) が線形同型であることを示した。
また
\[
I([\mu])=1
\]
であるから、\([\mu]\) は基底を与える。
\end{proof}

\begin{remark}
\(H^2(T^2)\) は、トーラス全体という閉じた $2$ 次元の面に対する総量
\[
\int_{T^2}\sigma
\]
で決まる。したがって、ここで検出されているのは「真ん中の穴」ではなく、
閉じた $2$ 次元多様体全体そのものである。
\end{remark}

\section{まとめ}

以上より
\[
H^0_{\mathrm{dR}}(T^2)\cong \R,\qquad
H^1_{\mathrm{dR}}(T^2)\cong \R^2,\qquad
H^2_{\mathrm{dR}}(T^2)\cong \R,
\]
また \(\dim T^2=2\) なので
\[
H^k_{\mathrm{dR}}(T^2)=0\qquad (k\ge 3)
\]
である。

基底は
\[
H^0:\ [1],\qquad
H^1:\ [\alpha],[\beta],\qquad
H^2:\ [\mu]
\]
で与えられる。ここで \(\alpha,\beta,\mu\) は
\[
\pi^*\alpha=dx,\qquad \pi^*\beta=dy,\qquad \pi^*\mu=dx\wedge dy
\]
を満たす \(T^2\) 上の形式である。

したがってコホモロジー環は
\[
H^*_{\mathrm{dR}}(T^2)\cong \Lambda_{\R}([\alpha],[\beta])
\]
と書ける。すなわち次数 $1$ の2つの生成元 \([\alpha],[\beta]\) を持つ外積代数である。

\section*{補論：縮められる loop に沿う閉 \(1\)-形式の積分}

\begin{prop}
\(\omega\in\Omega^1(T^2)\) を閉 \(1\)-形式とし、
\(\gamma:S^1\to T^2\) を loop とする。

さらに \(\gamma\) が縮められる、すなわち
ある滑らかな写像
\[
F:D^2\to T^2
\]
が存在して、その境界で
\[
F|_{\partial D^2}=\gamma
\]
となるとする。このとき
\[
\int_\gamma \omega=0
\]
が成り立つ。
\end{prop}

\begin{proof}
Stokes の定理より
\[
\int_\gamma \omega
=
\int_{\partial D^2}F^*\omega
=
\int_{D^2} d(F^*\omega)
\]
である。

pullback と外微分は可換だから
\[
d(F^*\omega)=F^*(d\omega)
\]
が成り立つ。しかも \(\omega\) は閉形式なので
\[
d\omega=0
\]
である。したがって
\[
\int_{D^2} d(F^*\omega)
=
\int_{D^2} F^*(d\omega)
=
\int_{D^2} F^*(0)
=
0.
\]
よって
\[
\int_\gamma \omega=0
\]
を得る。
\end{proof}

\begin{remark}
この命題はトーラスに特有のものではなく、任意の滑らかな多様体で成り立つ。
本稿で言う「穴に引っかかっていない loop」とは、この命題の仮定のように
\(D^2\) で張れる loop のことである。
\end{remark}

\section*{補足：\(g\) の周期性の別証明}

\begin{remark}
\(H^1_{\mathrm{dR}}(T^2)\) の計算において、\(\Phi\) の単射性を示す際、
\[
dg=\pi^*\omega
\]
を満たす \(g\in C^\infty(\R^2)\) の \(\Z^2\)-周期性は、座標計算の代わりに
整数平行移動
\[
\tau_{m,n}:\R^2\to\R^2,\qquad \tau_{m,n}(x,y)=(x+m,y+n)
\qquad ((m,n)\in \Z^2)
\]
を用いて、より概念的に示すこともできる。

\medskip

実際、\(\pi\circ \tau_{m,n}=\pi\) であるから
\[
d(g\circ \tau_{m,n}-g)
=
\tau_{m,n}^*(dg)-dg
=
\tau_{m,n}^*(\pi^*\omega)-\pi^*\omega
=
(\pi\circ \tau_{m,n})^*\omega-\pi^*\omega
=
\pi^*\omega-\pi^*\omega
=
0.
\]
したがって \(\R^2\) の連結性より
\[
g\circ \tau_{m,n}-g
\]
は定数関数である。これを \(c_{m,n}\in \R\) と書く。すなわち
\[
g(x+m,y+n)-g(x,y)=c_{m,n}
\qquad ((m,n)\in \Z^2).
\tag{A}
\]

\medskip

とくに \((m,n)=(1,0)\), \((0,1)\) の場合、\((A)\) の左辺は定数なので、
\((x,y)=(0,0)\) を代入して
\[
c_{1,0}=g(1,0)-g(0,0),\qquad c_{0,1}=g(0,1)-g(0,0)
\]
を得る。

ここで
\[
\widetilde{\gamma}_x(t)=(t,0),\qquad
\widetilde{\gamma}_y(t)=(0,t)
\qquad (0\le t\le 1)
\]
とすると、
\[
\pi\circ \widetilde{\gamma}_x=\gamma_x,\qquad
\pi\circ \widetilde{\gamma}_y=\gamma_y
\]
であるから
\[
c_{1,0}
=
g(1,0)-g(0,0)
=
\int_{\widetilde{\gamma}_x} dg
=
\int_{\widetilde{\gamma}_x} \pi^*\omega
=
\int_{\gamma_x}\omega,
\]
\[
c_{0,1}
=
g(0,1)-g(0,0)
=
\int_{\widetilde{\gamma}_y} dg
=
\int_{\widetilde{\gamma}_y} \pi^*\omega
=
\int_{\gamma_y}\omega.
\]
したがって、\(\Phi([\omega])=(0,0)\) の仮定
\[
\int_{\gamma_x}\omega=0,\qquad \int_{\gamma_y}\omega=0
\]
から
\[
c_{1,0}=0,\qquad c_{0,1}=0
\]
が従う。

\medskip

さらに \((A)\) から \(c_{m,n}\) は加法的である。実際、
\[
c_{m+m',\,n+n'}
=
g(x+m+m',y+n+n')-g(x,y)
\]
\[
=
\bigl(g(x+m+m',y+n+n')-g(x+m',y+n')\bigr)
+
\bigl(g(x+m',y+n')-g(x,y)\bigr)
=
c_{m,n}+c_{m',n'}
\]
である。よって
\[
c_{m,n}=m\,c_{1,0}+n\,c_{0,1}=0
\qquad ((m,n)\in \Z^2).
\]
したがって
\[
g(x+m,y+n)=g(x,y)
\qquad ((m,n)\in \Z^2)
\]
が成り立ち、\(g\) は \(\Z^2\)-周期的である。

ゆえに \(g\) は \(T^2\) 上の関数に降り、\(\omega\) が exact であることが従う。

\medskip

この議論は、\(g\) の周期性が「deck transformation によるずれは定数であり、その定数は
閉 \(1\)-形式 \(\omega\) の period に等しい」という事実から従うことを示している。
\end{remark}

\section*{付録：本稿で用いた被覆空間に関する事実}

この節では、\(T^2=\R^2/\Z^2\) の de Rham コホモロジーの計算で用いた
被覆空間に関する補足事項をまとめる。

\subsection*{基本設定}

商写像を
\[
\pi:\R^2\to T^2=\R^2/\Z^2,\qquad \pi(x,y)=[(x,y)]
\]
とする。ここで \(\Z^2\) は整数平行移動
\[
\tau_{m,n}:\R^2\to\R^2,\qquad \tau_{m,n}(x,y)=(x+m,y+n)
\qquad ((m,n)\in\Z^2)
\]
によって \(\R^2\) に作用する。

\begin{prop}[path lifting]
\(\gamma:[0,1]\to T^2\) を連続 path とし、\(p_0\in\R^2\) が
\[
\pi(p_0)=\gamma(0)
\]
を満たすとする。このとき、\(\gamma\) を始点 \(p_0\) から持ち上げる lift
\[
\widetilde{\gamma}:[0,1]\to\R^2
\]
で
\[
\pi\circ\widetilde{\gamma}=\gamma,\qquad \widetilde{\gamma}(0)=p_0
\]
を満たすものがただ一つ存在する。
\end{prop}

\begin{remark}
本稿で用いた
\[
\gamma_x(t)=\pi(t,0),\qquad \gamma_y(t)=\pi(0,t)
\]
の lift
\[
\widetilde{\gamma}_x(t)=(t,0),\qquad \widetilde{\gamma}_y(t)=(0,t)
\]
はこの一般事実の具体例である。
\end{remark}

\begin{prop}[lift を用いた pullback の計算]
\(\eta\in\Omega^k(T^2)\)、\(\gamma:[0,1]\to T^2\) を path、
\(\widetilde{\gamma}:[0,1]\to\R^2\) をその lift とする。このとき
\[
\gamma=\pi\circ\widetilde{\gamma}
\]
であるから
\[
\gamma^*\eta
=
(\pi\circ\widetilde{\gamma})^*\eta
=
\widetilde{\gamma}^{\,*}(\pi^*\eta)
\]
が成り立つ。
\end{prop}

\begin{proof}
pullback の合成則
\[
(f\circ g)^*=g^*\circ f^*
\]
を用いればよい。
\end{proof}

\begin{remark}
したがって \(T^2\) 上の形式の path に沿う積分は、\(\R^2\) 上の lift に沿う積分に直して計算できる。
本稿で
\[
\int_{\gamma_x}\alpha,\quad \int_{\gamma_x}\beta,\quad
\int_{\gamma_y}\alpha,\quad \int_{\gamma_y}\beta
\]
を計算したのはこの事実による。
\end{remark}

\begin{prop}[deck transformation]
各整数平行移動
\[
\tau_{m,n}(x,y)=(x+m,y+n)
\qquad ((m,n)\in\Z^2)
\]
は \(\pi\) の deck transformation である。すなわち
\[
\pi\circ\tau_{m,n}=\pi
\]
が成り立つ。さらに \(\pi\) の deck transformation はすべてこの形で与えられる。
\end{prop}

\begin{remark}
とくに、ある path \(\gamma\) の lift \(\widetilde{\gamma}\) が一つ与えられると、
他の任意の lift は
\[
\tau_{m,n}\circ \widetilde{\gamma}
\qquad ((m,n)\in\Z^2)
\]
の形に書ける。
\end{remark}

\begin{prop}[deck transformation と原始関数のずれ]
\(\omega\in\Omega^1(T^2)\) を閉 \(1\)-形式とし、
\[
dg=\pi^*\omega
\]
を満たす \(g\in C^\infty(\R^2)\) をとる。このとき任意の \((m,n)\in\Z^2\) に対して
\[
g\circ\tau_{m,n}-g
\]
は定数関数である。
\end{prop}

\begin{proof}
\[
d(g\circ\tau_{m,n}-g)
=
\tau_{m,n}^*(dg)-dg
=
\tau_{m,n}^*(\pi^*\omega)-\pi^*\omega
=
(\pi\circ\tau_{m,n})^*\omega-\pi^*\omega
=
0.
\]
\(\R^2\) は連結だから、\(g\circ\tau_{m,n}-g\) は定数である。
\end{proof}

\begin{remark}
この定数は \(\omega\) の period と解釈できる。
したがって、基本ループに沿う period がすべて \(0\) なら、
\(g\) は \(\Z^2\)-不変、すなわち \(T^2\) 上の関数に降りる。
これは \(H^1_{\mathrm{dR}}(T^2)\) の計算における \(\Phi\) の単射性の別証明を与える。
\end{remark}

\end{document}
